\documentclass[14pt,a4paper]{extarticle}
\usepackage[T2A]{fontenc}
\usepackage[utf8]{inputenc}
\usepackage[english,russian]{babel}
\usepackage{lmodern}
\usepackage{geometry}
\usepackage{array}
\usepackage{tabularx}
\usepackage{longtable}
\usepackage{ltablex}
\usepackage{booktabs}
\usepackage{hyperref}
\usepackage{indentfirst}
\usepackage{microtype}
\usepackage{graphicx}
\usepackage{setspace}
\usepackage{listings}
\usepackage{graphicx}
\usepackage{caption}
\usepackage{float}
\usepackage{tocloft}
\usepackage{titlesec}
\usepackage{newtxtext,newtxmath}

\titleformat{\section}        % форматируем \section
  {\normalfont\bfseries}      % стиль: обычный + жирный
  {\thesection.}              % номер секции с точкой
  {2mm}                       % расстояние между номером и текстом
  {}                           % перед заголовком можно добавить что-то
\titlespacing*{\section}{1.5cm}{*3}{*1}  % {отступ слева}{отступ сверху}{отступ снизу}

\captionsetup{justification=centering}
\renewcommand{\ttdefault}{cmtt}
\keepXColumns
\setlength{\parindent}{3em}
\newgeometry{left=2.5cm,right=2.5cm,top=3cm,bottom=2.5cm}

\renewcommand{\contentsname}{\hspace*{4em}{\smallСодержание}}
\setlength{\cftbeforesecskip}{0.5em}
\renewcommand{\cftsecleader}{\cftdotfill{\cftdotsep}}
\renewcommand{\cftsecaftersnum}{.\hspace{2mm}}

\title{Отчёт по "Организация работы ИТ-службы"}
\author{Выполнил студент группы М-РИТ-25-1, Станиславчук Сергей}

\begin{document}

\setstretch{1.5}
\noindent
\begin{minipage}[t]{0.01\textwidth}
    \includegraphics{~/programmer/++Programmer/5_1/LSTU.jpg}
\end{minipage}%
\begin{minipage}[c]{1.1\linewidth}
    \centering
    {\small
    МИНИСТЕРСТВО НАУКИ И ВЫСШЕГО ОБРАЗОВАНИЯ РФ \\[1em]
    ФЕДЕРАЛЬНОЕ ГОСУДАРСТВЕННОЕ БЮДЖЕТНОЕ \\
    ОБРАЗОВАТЕЛЬНОЕ УЧРЕЖДЕНИЕ ВЫСШЕГО ОБРАЗОВАНИЯ \\
    «ЛИПЕЦКИЙ ГОСУДАРСТВЕННЫЙ ТЕХНИЧЕСКИЙ УНИВЕРСИТЕТ»
    }
\end{minipage}%

\vspace{3em}

\noindent
\hspace{1em}%
\begin{minipage}[t]{0.8\textwidth}
\textbf{Институт} \hspace{2em} компьютерных наук \\
\textbf{Кафедра} \hspace{2em} автоматизированных систем управления
\end{minipage}


\vspace{6em}

\begin{center}
\textbf{Практическая работа №2} \\
по дисциплине \\ <<Организация работы ИТ-службы>>
\end{center}

\vspace{6em}

\noindent
\begin{tabular*}{\textwidth}{@{\extracolsep{\fill}} l c r }
    Студент М-РИТ-25-1 & \makebox[4cm]{\hrulefill} & Станиславчук С.М. \\[2em]
Руководитель & \\
К.т.н., доцент & \makebox[4cm]{\hrulefill} & Ведищев В.В. \\
\end{tabular*}

\vfill
\begin{center}
    Липецк, 2026
\end{center}

\thispagestyle{empty}

\newpage
\tableofcontents
\newpage

\section{Стратегическое управление ИТ-услугами}

Стратегическое управление ИТ-услугами направлено на определение целей предоставления услуги, её ценности для заказчиков и направлений дальнейшего развития. В случае VDS-хостинга основная цель заключается в предоставлении клиентам вычислительных ресурсов для размещения собственных информационных систем без необходимости приобретения и обслуживания физических серверов.

При формировании стратегии необходимо учитывать потребности заказчиков, развитие технологий виртуализации, действия конкурентов, стоимость инфраструктуры и требования информационной безопасности. Например, увеличение спроса на ресурсоёмкие приложения потребует расширения вычислительных мощностей и модернизации серверного оборудования.

Стратегическое управление также предполагает определение приоритетов развития сервиса. Поставщик должен поддерживать баланс между качеством услуги, её стоимостью, доступностью и рентабельностью. Результатом стратегического управления является определение направлений развития VDS-хостинга и требований к предоставляемой услуге.

\section{Управление портфелем услуг}

Управление портфелем услуг представляет собой процесс управления всеми услугами поставщика на протяжении их жизненного цикла. Он позволяет определять, какие услуги необходимо развивать, поддерживать, изменять или выводить из эксплуатации.

В случае VDS-хостинга портфель включает различные конфигурации виртуальных серверов, отличающиеся количеством виртуальных процессорных ядер, объёмом оперативной памяти, дисковым пространством и сетевыми характеристиками. Дополнительно могут предоставляться услуги хранения данных и другие возможности, связанные с эксплуатацией виртуальной инфраструктуры.

Управление портфелем позволяет оценивать востребованность отдельных предложений и эффективность использования ресурсов. Например, если спрос на конфигурации с большим объёмом оперативной памяти растёт, поставщик может расширить линейку соответствующих тарифов. Если определённая конфигурация становится невостребованной или перестаёт быть экономически целесообразной, её можно изменить или исключить из предложения.

Таким образом, управление портфелем помогает согласовывать состав услуг с потребностями рынка и стратегическими целями поставщика.

\section{Управление финансами для ИТ-услуг}

Управление финансами для ИТ-услуг обеспечивает планирование, учёт и контроль затрат, связанных с предоставлением услуги, а также оценку её экономической эффективности.

Для VDS-хостинга основные затраты включают приобретение и обслуживание физических серверов, систем хранения данных и сетевого оборудования, оплату размещения в дата-центрах, электроэнергию, сетевые услуги, лицензии на программное обеспечение и заработную плату технического персонала.

Поставщик должен соотносить стоимость предоставления услуги с доходами от клиентских тарифов. При определении цены необходимо учитывать не только непосредственные затраты на оборудование, но и расходы на техническую поддержку, обеспечение безопасности, резервирование и модернизацию инфраструктуры.

Например, увеличение количества клиентов повышает потенциальную выручку, однако одновременно увеличивает нагрузку на физические серверы и инфраструктуру. Если ресурсов становится недостаточно, поставщику требуется инвестировать в расширение мощностей. Поэтому финансовое управление должно обеспечивать рациональное распределение ресурсов и поддержание экономической устойчивости сервиса.

\section{Управление спросом}

Управление спросом направлено на выявление и прогнозирование потребности заказчиков в ИТ-услуге, а также на согласование этой потребности с возможностями поставщика.

Для VDS-хостинга спрос определяется количеством клиентов, выбираемыми конфигурациями серверов, требованиями к производительности и характером размещаемых приложений. Например, веб-сайты с небольшой посещаемостью могут работать на конфигурациях с ограниченными ресурсами, тогда как базы данных и высоконагруженные приложения требуют большего количества процессорных ресурсов, оперативной памяти и дисковой производительности.

Поставщик должен анализировать динамику заказов, загрузку физических серверов и распределение спроса между тарифами. На основании этих данных планируется закупка оборудования, расширение инфраструктуры и изменение доступных конфигураций.

Недостаточное прогнозирование спроса приводит к нехватке ресурсов и снижению качества обслуживания. Избыточное расширение инфраструктуры, напротив, приводит к неэффективному использованию капитала и увеличению расходов. Следовательно, задача управления спросом заключается в согласовании доступных мощностей с текущими и ожидаемыми потребностями клиентов.

\section{Управление взаимоотношениями с бизнесом}

Управление взаимоотношениями с бизнесом направлено на взаимодействие поставщика ИТ-услуги с заказчиками, выявление их потребностей и обеспечение соответствия услуги ожидаемым результатам.

В случае VDS-хостинга заказчиками выступают частные лица, разработчики, малые и средние предприятия, которым требуется инфраструктура для размещения приложений и информационных систем. Поставщик должен понимать требования каждой категории клиентов к производительности, доступности, безопасности и стоимости услуги.

Взаимодействие с заказчиками осуществляется через каналы технической поддержки, обращения по вопросам обслуживания, консультации и процедуры обработки жалоб. Например, если клиент сообщает о недостаточной производительности виртуального сервера, поставщик должен определить характер проблемы и предложить подходящее решение: изменение конфигурации, увеличение ресурсов или устранение неисправности инфраструктуры.

Важную роль играет согласование ожиданий заказчика с фактическими возможностями услуги. Для этого определяются условия предоставления сервиса, параметры качества и ответственность сторон. Обратная связь от клиентов позволяет выявлять недостатки обслуживания, корректировать предложения и определять направления улучшения.

\end{document}
